\begin{table*}[htbp]
\centering\small
\caption{Utilidade percebida.}
\label{tab:usuarios-utilidade}
\begin{tabular}{p{0.37\textwidth}cccp{0.18\textwidth}cc}
\hline
Recurso & n & Med. & Q1--Q3 & 4--5 & 1--2 & NC\\ \hline
Busca por localização geográfica & 32 & 4 & 4,0--5,0 & 26/32 (81,2\%) & 3 & 0\\
Demarcação de áreas agrícolas & 32 & 4 & 3,0--5,0 & 22/32 (68,8\%) & 2 & 0\\
Salvamento e gerenciamento de áreas & 24 & 4 & 4,0--4,0 & 19/24 (79,2\%) & 1 & 8\\
Análise de índices espectrais & 31 & 4 & 3,0--5,0 & 21/31 (67,7\%) & 5 & 1\\
Mapas temáticos com gradientes de cores & 31 & 5 & 4,0--5,0 & 25/31 (80,6\%) & 1 & 1\\
Consulta interativa de valores por pixel & 24 & 5 & 4,0--5,0 & 21/24 (87,5\%) & 0 & 8\\
Gráficos de evolução temporal & 30 & 4 & 3,0--5,0 & 21/30 (70,0\%) & 4 & 2\\
Detecção de anomalias & 30 & 4 & 4,0--4,8 & 25/30 (83,3\%) & 1 & 2\\
Informações de temperatura e precipitação & 30 & 5 & 4,0--5,0 & 24/30 (80,0\%) & 0 & 2\\
Geração de relatórios PDF & 28 & 4 & 4,0--5,0 & 23/28 (82,1\%) & 1 & 4\\
Exportação em Excel e CSV & 27 & 4 & 4,0--4,0 & 21/27 (77,8\%) & 1 & 5\\
\hline
\end{tabular}
\fonteTCC{dados da pesquisa (2026).}
\notaTCC{Med.: mediana; Q1--Q3: quartis; 4--5: útil/muito útil; 1--2: nada/pouco útil; NC: não conseguiu avaliar. Percentuais usam o n válido do item.}
\end{table*}
